\chapter{Einleitung} \label{ch:Einleitung}

\section{Aufgabe und Motivation}
\label{sec:AuM}
Dieser Versuch beschäftigt sich mit mit drei der Grundlegensten Bauteilen der Elektrotechnik; dem ohmischen Widerstand, dem Kondensator und der Spule. Je nach Anordnung und Kombination der Bauteile lassen sich verschiedene Eigenschaften nutzen. So ist das herausfiltern bestimmter Frequenzen möglich, auch eine Annährung an die Integration und die Differentation eines Signals lässt sich via jener bauteile möglich.

\section{Physikalische Grundlage}
\label{sec:PhyBasics}


\subsubsection*{Skizze des Versuchsaufbaus}
\label{sec:Skizze}

